\section{Exam day, 2026-12-05}

Everything in the month points at two sessions of three hours. This section is the protocol for those six hours, and the checklist for the pages that get handed in.

\subsection*{Before the morning session}

Sleep. Eat something with protein. Bring several pens, a watch that is not a phone, and water. Do not read problems on the way in. The reading habit from \secref{sec:reading}{reading a problem} is the whole preparation; nothing learned in the last hour will be available under pressure.

\subsection*{The first fifteen minutes}

Read all six problems before writing anything. Then, on scratch paper, one call line for each, with a rank from 1 to 6:

\begin{quote}
\texttt{A1: <technique>, <first assumption>, rank 1. 70}
\end{quote}

No times, because the box on the day is set by the rank. This is the one habit from the month that pays directly. Six call lines take fifteen minutes and produce a plan for the remaining 165, and they stop the most common failure, which is spending the first hour on A3 because its statement was interesting.

The ranking is not by slot. A1 is usually the easiest, but not always, and a B2 that matches a ladder you climbed last week outranks a B1 that does not. Rank by confidence that you can finish, not by how much you like the problem.

\subsection*{The remaining 165 minutes}

The top-ranked problem until it is solved and written up. Then the second. Then the third if there is time, which there usually is not.

A problem that is not falling after 45 minutes gets dropped for the next on the rank, without regret. The score distribution in Section~1 is the argument: a finished second problem is worth 10 and a half-finished first is worth 1, so the move from a stuck problem to a fresh one is a gain of about nine points in expectation, even if the fresh one is harder. Write ``dropped at <time>'' on the scratch call line so the evening ledger is honest.

If the top-ranked problem falls quickly, write it up fully before moving on. The temptation is to leave the write-up for later and chase the second problem while the first is fresh; the result is often two 1s instead of a 10 and a 1.

Fifteen minutes before the end, stop solving. Check what is written against the checklist below. A 9 recovered to a 10 in those fifteen minutes is worth more than any progress on a new problem.

\subsection*{Writing up}

Graders give 9 and 10 to arguments they can follow without goodwill, and 0 to ideas. The write-up is a separate skill from the solving, and this is the checklist.

\begin{enumerate}[nosep,leftmargin=*]
\item State what you are proving at the top. One sentence. If the problem says ``find all'', state the answer there too.
\item Define every symbol you introduce, at its first use. ``Let $d$ be the smallest prime factor of $n$'' before $d$ appears in a formula.
\item Make every case explicit, including the one that is obvious. If the argument splits on $n$ even or odd, both branches are on the page, even if one is a line.
\item The equality case of an inequality, the base case of an induction, the $n = 1$ case of a sum: write them. These are the gaps that turn 10 into 9.
\item If the answer is a value or a set, box it, and say why no other value works. A correct value with no uniqueness argument is a 1.
\item Any fact used from outside the problem is named: Fermat's little theorem, Rolle, the pigeonhole principle. Named facts can be checked; unnamed ones look like gaps.
\item A construction is checked against every condition in the statement, in writing. ``This works because'' followed by the verification, not followed by nothing.
\item Nothing in the final write-up refers to the scratch work. Cross out what is not part of the argument. A grader who has to decide which of two attempts is the real one will pick neither.
\item Read the whole solution once as a stranger would, looking for a step that uses something not yet shown.
\end{enumerate}

Write in full sentences. A line of symbols with no verb is not an argument. ``Since $p$ does not divide $a$, Fermat gives $a^{p-1} \equiv 1 \pmod p$'' is a sentence; the same with the words removed is a hope.

\subsection*{Between sessions}

Eat. Do not discuss the morning problems with anyone, and do not look anything up. Whatever happened in the morning has no bearing on the afternoon, and the afternoon starts with the same fifteen minutes of reading.

\subsection*{That evening}

Type the twelve call lines into \texttt{calls.md}, with the date and the domain \texttt{putnam}, exactly as they were written on paper, including the ranks and the ones marked dropped. The unattempted ones are \emph{moot}. Settle them all when the archive posts the solutions, usually within the week. Run \texttt{/profile} the weekend after; the \texttt{putnam} row under Domains, and its lines under Misses, are the evidence the month was for.

